\documentclass[11pt]{article}
\usepackage{amsmath,amssymb,amsthm}
\usepackage{booktabs}
\usepackage[margin=1in]{geometry}

\newtheorem{theorem}{Theorem}
\newtheorem{lemma}[theorem]{Lemma}
\newtheorem{corollary}[theorem]{Corollary}
\theoremstyle{definition}
\newtheorem{conjecture}[theorem]{Conjecture}

\title{A Disproof of Conjecture 00000000155:\\
There Are More Bell Primes Than $B_2$ and $B_3$}
\author{earthking11}
\date{September 13, 2026}

\begin{document}
\maketitle

\begin{abstract}
Conjecture 00000000155 asserts that the only Bell primes are $B_2 = 2$ and
$B_3 = 5$. We refute it: in the same $0$-indexed convention used by the
conjecture file, $B_7 = 877$ and $B_{13} = 27644437$ are both prime, and
neither equals $2$ or $5$. The claim is therefore false. The computation is
formalized in core Lean~4 (\texttt{import Std}, no Mathlib) and reproduced
with an independent Python script.
\end{abstract}

\section{The Bell numbers}

The Bell number $B_n$ counts the partitions of an $n$-element set, with
$B_0 = B_1 = 1$. The conjecture file fixes the convention by writing
$B_2 = 2$ and $B_3 = 5$; this is the standard $0$-indexed sequence
OEIS A000110. The first values, computed with the Bell triangle
(Aitken's array),
\[
a_{0,0} = 1,\qquad
a_{n,0} = a_{n-1,n-1},\qquad
a_{n,k} = a_{n,k-1} + a_{n-1,k-1},
\]
and $B_n = a_{n,0}$, are:

\begin{center}
\begin{tabular}{cc}
\toprule
$n$ & $B_n$ \\
\midrule
$0$  & $1$ \\
$1$  & $1$ \\
$2$  & $2$ \quad (prime) \\
$3$  & $5$ \quad (prime) \\
$4$  & $15$ \\
$5$  & $52$ \\
$6$  & $203$ \\
$7$  & $877$ \quad (prime) \\
$8$  & $4140$ \\
$9$  & $21147$ \\
$10$ & $115975$ \\
$11$ & $678570$ \\
$12$ & $4213597$ \\
$13$ & $27644437$ \quad (prime) \\
\bottomrule
\end{tabular}
\end{center}

\section{Primality certificates}

\begin{lemma}
$877$ is prime.
\end{lemma}

\begin{proof}
Since $\sqrt{877} < 30$, it suffices to rule out the ten primes
$2,3,5,7,11,13,17,19,23,29$. Direct division gives the nonzero remainders
\[
\begin{aligned}
877 \bmod 2 = 1, &\qquad
877 \bmod 3 = 1, &\qquad
877 \bmod 5 = 2,\\
877 \bmod 7 = 2, &\qquad
877 \bmod 11 = 8, &\qquad
877 \bmod 13 = 6,\\
877 \bmod 17 = 10, &\qquad
877 \bmod 19 = 3, &\qquad
877 \bmod 23 = 3,\\
877 \bmod 29 = 7. &&
\end{aligned}
\]
No prime $p \le 29$ divides $877$, so $877$ is prime.
\end{proof}

\begin{lemma}
$27644437$ is prime.
\end{lemma}

\begin{proof}
Here $\sqrt{27644437} < 5258$, since $5257^2 = 27636049 \le 27644437 <
27646564 = 5258^2$.  It therefore suffices to check that no odd
$d$ with $3 \le d \le 5257$ divides $27644437$ (the even cases are excluded
because the number is odd).  An exhaustive trial division, performed exactly
in integer arithmetic, finds no divisor; equivalently, the Lean function
\texttt{isPrimeB} returns \texttt{true} on this input, and the Python
reproduction script confirms the same.  The formal trial-division loop stops
as soon as $d^2 > n$, so it examines only $d = 3, 5, \dots, 5257$, a finite
check of $2628$ candidates.
\end{proof}

\section{Refutation}

\begin{theorem}
The only Bell primes are not $B_2 = 2$ and $B_3 = 5$.
\end{theorem}

\begin{proof}
By the table, $B_7 = 877$ and $B_{13} = 27644437$, and by the two lemmas
both are prime. Since $877 \notin \{2, 5\}$ and
$27644437 \notin \{2, 5\}$, the set of Bell primes contains at least
$\{2, 5, 877, 27644437\}$, which is strictly larger than $\{2, 5\}$.
Hence the conjecture is false. Formally, in \texttt{Main.lean} we prove
\[
\texttt{bell } 2 = 2 \;\wedge\; \texttt{bell } 3 = 5 \;\wedge\;
\texttt{isPrimeB }(\texttt{bell } 7) = \texttt{true} \;\wedge\;
\texttt{isPrimeB }(\texttt{bell } 13) = \texttt{true},
\]
together with $\texttt{bell } 7 \ne \texttt{bell } 2$ and
$\texttt{bell } 7 \ne \texttt{bell } 3$, all by kernel reduction with
\texttt{decide} (no \texttt{sorry}, no \texttt{axiom}, no
\texttt{native\_decide}).
\end{proof}

\begin{conjecture}[00000000155, refuted]
The only Bell primes are $B_2 = 2$ and $B_3 = 5$.
\end{conjecture}

\begin{corollary}
The list of Bell primes begins $2, 5, 877, 27644437, \dots$; in particular
the conjecture is incomplete.  We verify the range $0 \le n \le 20$ here;
larger Bell primes are reported in the literature, but are not needed for
the refutation.
\end{corollary}

\end{document}
